%% copyright --- the page after the title page.
%%
%% \bookedition and \bookyear are pinned in preamble.tex and \pinnedon is a
%% user-visible string in lang/<lang>.tex. None of the three is written out by
%% hand here.

\thispagestyle{empty}
\vspace*{\fill}

{\small\color{mgrey}\raggedright\setlength{\parindent}{0pt}%
\emph{Chaos od zera. Jak matematyka opisuje świat i gdzie przestaje go
przewidywać}. Wydanie \bookedition.

\medskip
Copyright \textcopyright{} \bookyear{} Konrad Cinkusz. Wszelkie prawa
zastrzeżone, z wyjątkiem udzielonych w poniższych licencjach.

\medskip
Tekst jest udostępniony na licencji \proselicence{}. Kod \dash{} każdy plik
w katalogu \texttt{code/} i każdy listing, który książka z niego drukuje
\dash{} na licencji MIT.

\medskip
Wydanie własne autora. Książka nie ma wydawcy, dystrybutora ani numeru
ISBN.

\medskip
Wydana w dwóch językach, polskim i angielskim, z jednego repozytorium.
Dostępna także po angielsku jako \emph{Chaos from Zero}.

\medskip
Wersje ustalone: \pinnedon. Każdą liczbę w tej książce, której nie da się
policzyć w pamięci, wyprodukował program z repozytorium, a proces budowania
uruchamia wszystkie te programy ponownie przy każdej zmianie i odmawia
zbudowania książki, jeśli choć jedna wydrukowana cyfra przestała się zgadzać.

\medskip
Ta książka nie należy do żadnej serii poradników dla początkujących,
niezależnie od tego, jak nazywa się repozytorium, w którym powstała.
\par}

\clearpage
